\documentclass[runningheads]{llncs}
\usepackage{graphicx}
\usepackage{amsmath}
\usepackage{booktabs}
\usepackage{hyperref}
\usepackage{float}

\begin{document}

\title{Building an End-to-End Multi-Climate Temperature Prediction System: A Modular Pipeline and Astronomical Feature Approach}

\author{Do Nguyen Huu Phuc\inst{1}\orcidID{0009-0001-1761-3423} \and
Ho Truong Thinh\inst{1}\orcidID{0009-0006-3125-665X} \and
Dinh Tran Khanh Nam\inst{1}\orcidID{0009-0000-7903-362X} \and
Pham Truong Minh Nhat\inst{1}\orcidID{0009-0005-7363-3591}}

\institute{FPT University, Ho Chi Minh City, Viet Nam \\
\email{\{do.huuphuc2k5, hotruongthinh2000, namlamborghini2404, minhnhat9933777\}@gmail.com} \\
\textbf{Instructor:} Le Vo Minh Thu}

\maketitle

\begin{abstract}
Traditional numerical weather prediction (NWP) models demand immense computational resources, whereas standard machine learning models often suffer from geographical overfitting, failing to generalize across diverse regions. This paper proposes a robust, end-to-end machine learning pipeline specifically engineered to predict temperatures across five distinct European climate zones. The proposed system features a meticulously designed four-layer modular architecture, integrating automated API data ingestion, strict SQL-based anomaly detection, and the parallel training of 15 specialized models. To overcome the limitations of raw meteorological data, we introduce an advanced domain-knowledge feature engineering framework. This includes the trigonometric calculation of solar declination for astronomical daylight derivation, temporal lag generation to capture atmospheric thermal inertia, and cyclical encoding for time-series continuity. Experimental evaluations demonstrate that Support Vector Regression (SVR) achieves superior predictive performance, yielding a Coefficient of Determination ($R^2$) of approximately 0.949 and a Mean Absolute Error (MAE) of 1.32$^\circ$C. Furthermore, utilizing model persistence, the deployed Streamlit web application guarantees real-time inference with latency under 100 milliseconds, establishing its viability for commercial applications such as energy optimization and logistics planning.

\keywords{Multi-Climate Prediction \and Support Vector Regression \and Astronomical Features \and Cyclical Encoding \and Modular Architecture.}
\end{abstract}

\section{Introduction}

\subsection{Context and Problem Statement}
Accurate meteorological forecasting remains a critical component for disaster management, agricultural planning, and energy optimization. Currently, state-of-the-art predictions rely heavily on Numerical Weather Prediction (NWP) systems, which simulate complex atmospheric physics but require supercomputers to operate. Conversely, localized data-driven machine learning models provide rapid inferences but frequently ignore core physical phenomena, leading to severe spatial overfitting. The primary motivation of this research is to bridge this gap by constructing a lightweight, mathematically grounded, and highly generalizable forecasting system capable of adapting to five distinctly classified Köppen climate zones in Europe: Oceanic (Amsterdam), Continental (Berlin), Mediterranean (Athens), Nordic (Stockholm), and Alpine (Zurich).

\subsection{Main Contributions}
The core contributions of this comprehensive study are manifold:
\begin{itemize}
    \item Formulating a fully automated, end-to-end modular architecture that dynamically processes raw API data into deployable predictive insights.
    \item Implementing an advanced feature engineering pipeline grounded in astronomical physics (solar declination) and time-series mathematics (cyclical encoding).
    \item Conducting a rigorous comparative analysis among five distinct machine learning algorithms to establish a robust baseline for multi-target temperature forecasting.
\end{itemize}

\section{System Architecture}
To achieve high scalability and ensure a sterile data environment, the system strictly adheres to a four-layer modular architecture. This separation of concerns prevents processing bottlenecks during real-time inference.

\begin{enumerate}
    \item \textbf{Data Ingestion Layer:} Automates the extraction process from the Open-Meteo API. It utilizes advanced temporal resampling techniques to uniformly merge high-frequency hourly metrics with daily temperature records.
    \item \textbf{Data Validation and EDA Layer:} Instead of relying solely on standard dataframes, this layer integrates an in-memory SQLite3 database. By executing complex SQL queries, the system efficiently isolates statistical anomalies and extreme weather events, ensuring the integrity of the training dataset.
    \item \textbf{Modeling Layer:} Employs a multi-model, multi-target paradigm. Five distinct algorithms are trained in parallel across three independent target variables (TempMax, TempMean, TempMin), resulting in an ensemble of 15 specialized models.
    \item \textbf{Presentation Layer:} The final phase involves deploying the optimized models into an interactive interface.
\end{enumerate}

\begin{figure}[H]
    \centering
    \includegraphics[width=\textwidth]{workflow_diagram.png}
    \caption{End-to-End Workflow illustrating the transition from raw API ingestion, through SQL-based validation and mathematical extraction, to model training.}
    \label{fig:workflow}
\end{figure}

\section{Application Implementation}
The theoretical models were successfully transitioned into a functional, user-centric software product using the Streamlit framework. A critical engineering challenge in deployment is minimizing response latency. To address this, the system leverages model persistence—serializing the fully trained models into binary `.pkl` formats. 

When a user inputs target parameters via the frontend interface, the backend instantaneously loads these serialized components, bypassing any computational overhead associated with retraining. This architectural decision enables the application to deliver highly accurate predictions with a latency of less than 100 milliseconds.

\begin{figure}[H]
    \centering
    \includegraphics[width=\textwidth]{system_architecture.png}
    \caption{Interactive deployment architecture demonstrating the real-time inference mechanism between the user inputs and the serialized backend models.}
    \label{fig:sys_arch}
\end{figure}

\section{Methodology and Feature Engineering}

\subsection{Data Collection and Processing}
The dataset encompasses comprehensive meteorological records for the five representative European cities. To rigorously evaluate the models and prevent the critical flaw of temporal data leakage, we discarded standard random splitting. Instead, a strict Chronological Split was enforced, ensuring that the models are trained exclusively on historical sequences and validated only on subsequent, unseen future data.

\begin{figure}[H]
    \centering
    \includegraphics[width=\textwidth]{data_distribution.png}
    \caption{Analysis of dataset balance across the five cities (left) and the target temperature range distribution demonstrating the study's broad scope (right).}
    \label{fig:data_dist}
\end{figure}

\subsection{Astronomical Feature Engineering}
A fundamental limitation of utilizing raw dates in machine learning is the algorithm's inability to comprehend geographic variations in solar radiation. Since temperature is intrinsically driven by the sun, we mathematically engineered an astronomical daylight feature based on latitudinal coordinates. 

First, the solar declination angle ($\delta$) is calculated to determine the Earth's axial tilt relative to the sun on any given day of the year ($N$):
\begin{equation}
    \delta = 23.45^\circ \sin\left[\frac{360}{365} (N - 81)\right]
\end{equation}
Subsequently, the total daylight hours ($D$) for a specific latitude ($\phi$) is derived using spherical trigonometry:
\begin{equation}
    D = \frac{24}{\pi} \arccos\left(-\tan(\phi) \tan(\delta)\right)
\end{equation}

\begin{figure}[H]
    \centering
    \includegraphics[width=\textwidth]{daylight_hours_variation.png}
    \caption{Astronomical daylight hours variation by latitude in 2023. Note: Due to the latitudinal proximity of Amsterdam and Berlin, distinct linestyles are applied to differentiate their overlapping trajectories.}
    \label{fig:daylight}
\end{figure}

\subsection{Feature Correlation and Autoregression}
Atmospheric temperature does not fluctuate randomly; it possesses a "thermal inertia," meaning today's climate is heavily influenced by the preceding days. To capture this autoregressive property, we engineered temporal lag variables (e.g., \texttt{TempLag1}, \texttt{TempLag3}). As validated by the correlation heatmap, these artificially engineered features, alongside the astronomical daylight hours, exhibit an exceptionally strong positive correlation with the target temperature, proving their high predictive value.

\begin{figure}[H]
    \centering
    \includegraphics[width=0.8\textwidth]{correlation_matrix.png}
    \caption{Correlation matrix confirming the high predictive power of temporal lags and daylight hours against the target temperature.}
    \label{fig:correlation}
\end{figure}

\section{Experimental Results and Evaluation}

\subsection{Model Comparison}
The experimental phase evaluated five distinct algorithms: Ridge Regression, K-Nearest Neighbors (KNN), Random Forest (RF), Extreme Gradient Boosting (XGBoost), and Support Vector Regression (SVR). Hyperparameter optimization was conducted to ensure a fair baseline. The results unequivocally demonstrate that SVR and XGBoost significantly outperform distance-based and linear models. 

Ultimately, Support Vector Regression (SVR) emerged as the optimal model, achieving an $R^2 \approx 0.949$ and an MAE of approximately 1.32$^\circ$C. SVR's superiority in this context can be attributed to its utilization of the kernel trick, which effectively maps the complex, non-linear climate features into a high-dimensional space, thereby maximizing the margin of tolerance for prediction errors across the diverse weather zones.

\begin{figure}[H]
    \centering
    \includegraphics[width=0.9\textwidth]{model_comparison.png}
    \caption{Performance leaderboard ($R^2$ Score) among the five machine learning algorithms, highlighting the top-tier accuracy of the SVR and XGBoost models.}
    \label{fig:model_comp}
\end{figure}

\subsection{Reliability Assessment}
To visually assess the degree of error, the actual ground-truth temperatures were plotted against the SVR predictions. The resulting scatter plot indicates that the data points corresponding to all five cities align tightly along the ideal prediction line ($y=x$). This confirms that the model is not only highly accurate but also geographically unbiased, successfully generalizing across diverse climates without skewing towards any specific region.

\begin{figure}[H]
    \centering
    \includegraphics[width=0.8\textwidth]{actual_vs_predicted.png}
    \caption{Correlation between actual and predicted temperature values, demonstrating high reliability and minimal regional bias.}
    \label{fig:actual_predicted}
\end{figure}

\subsection{Seasonal Error Analysis}
While the overall metric is exceptional, a granular temporal analysis of the Root Mean Square Error (RMSE) reveals minor seasonal fluctuations. The model maintains highly stable and precise predictions throughout the spring, summer, and autumn months. However, the error margin increases slightly during the winter season (December, January, and February). This variance is scientifically expected, as winter in Europe is characterized by volatile atmospheric pressure changes and unpredictable extreme weather anomalies (such as sudden blizzards or polar vortexes) which inherently complicate forecasting.

\begin{figure}[H]
    \centering
    \includegraphics[width=0.85\textwidth]{seasonal_performance.png}
    \caption{Average monthly RMSE analysis demonstrating consistent stability with expected minor error increments during extreme winter months.}
    \label{fig:seasonal_error}
\end{figure}

\section{Conclusion}
This research successfully bridges the gap between atmospheric physics and modern data science by deploying a robust, multi-climate temperature prediction system. Through the strategic implementation of a modular architecture, the system guarantees a clean data pipeline and ultra-fast inference. The empirical evidence strongly highlights the necessity of domain-specific feature engineering; the integration of astronomical daylight calculations and temporal lags provided the SVR model with the contextual depth required to achieve an $R^2 \approx 0.949$ and an MAE of 1.32$^\circ$C.

\begin{figure}[H]
    \centering
    \includegraphics[width=\textwidth]{tech_stack.png}
    \caption{Technology stack summarizing the integrated data, modeling, mathematical, and presentation layers utilized throughout the project.}
    \label{fig:tech_stack}
\end{figure}

The resulting application is not merely an academic exercise but a viable commercial tool ready for integration into energy management and logistics systems. Future development will focus on incorporating Deep Learning architectures, such as Long Short-Term Memory (LSTM) networks, to extend the temporal forecasting horizon and capturing additional variables like atmospheric pressure to predict precipitation.

\begin{thebibliography}{8}
\bibitem{ref_openmeteo}
Open-Meteo: Open-source Weather API. \url{https://open-meteo.com/} (2024).

\bibitem{ref_xgboost}
Chen, T., Guestrin, C.: XGBoost: A Scalable Tree Boosting System. In: Proceedings of the 22nd ACM SIGKDD International Conference on Knowledge Discovery and Data Mining, pp. 785--794 (2016).

\bibitem{ref_svr}
Smola, A.J., Schölkopf, B.: A tutorial on support vector regression. Statistics and Computing 14(3), 199--222 (2004).

\bibitem{ref_time}
Hyndman, R.J., Athanasopoulos, G.: Forecasting: Principles and Practice, 3rd edition. OTexts: Melbourne, Australia (2021).

\bibitem{ref_streamlit}
Streamlit: The fastest way to build and share data apps. \url{https://streamlit.io/} (2024).
\end{thebibliography}

\end{document}